\documentclass[10pt,letterpaper]{article}

\usepackage{fontspec}
\setmainfont{TeX Gyre Termes}
\setsansfont{TeX Gyre Heros}
\setmonofont{DejaVu Sans Mono}[Scale=0.82]

\usepackage[margin=0.82in,headheight=15pt]{geometry}
\usepackage{microtype}
\usepackage{amsmath}
\usepackage{amssymb}
\usepackage{booktabs}
\usepackage{array}
\usepackage{tabularx}
\usepackage{longtable}
\usepackage{adjustbox}
\usepackage{enumitem}
\usepackage{graphicx}
\usepackage{xcolor}
\usepackage[most]{tcolorbox}
\usepackage{titlesec}
\usepackage{fancyhdr}
\usepackage{siunitx}
\usepackage{pgfplots}
\usepackage{fancyvrb}
\usepackage{csquotes}
\usepackage[
  backend=biber,
  style=numeric,
  sorting=none,
  maxbibnames=9,
  url=true,
  doi=true
]{biblatex}
\usepackage[
  colorlinks=true,
  linkcolor=blue!45!black,
  citecolor=blue!45!black,
  urlcolor=blue!55!black,
  pdfauthor={CompactInVacuum Polarimeter Project},
  pdftitle={CompactInVacuum SiPM Detector, Vacuum Chamber, and China Procurement Engineering Report}
]{hyperref}
\usepackage{bookmark}
\usepackage[nameinlink,noabbrev]{cleveref}

\addbibresource{references.bib}
\pgfplotsset{compat=1.18}
\sisetup{
  detect-all,
  group-separator={,},
  group-minimum-digits=4,
  range-phrase={--},
  range-units=single
}

\definecolor{ReportBlue}{HTML}{2E74B5}
\definecolor{ReportDarkBlue}{HTML}{1F4D78}
\definecolor{ReportText}{HTML}{202124}
\definecolor{ReportMuted}{HTML}{5F6368}
\definecolor{ReportGreen}{HTML}{E8F3EC}
\definecolor{ReportAmber}{HTML}{FFF4D6}
\definecolor{ReportRed}{HTML}{FCE8E6}
\definecolor{ReportLight}{HTML}{F2F4F7}

\titleformat{\section}
  {\sffamily\bfseries\Large\color{ReportBlue}}
  {\thesection.}{0.55em}{}
\titleformat{\subsection}
  {\sffamily\bfseries\large\color{ReportBlue}}
  {\thesubsection}{0.55em}{}
\titleformat{\subsubsection}
  {\sffamily\bfseries\normalsize\color{ReportDarkBlue}}
  {\thesubsubsection}{0.55em}{}
\titlespacing*{\section}{0pt}{1.6em}{0.7em}
\titlespacing*{\subsection}{0pt}{1.25em}{0.45em}
\titlespacing*{\subsubsection}{0pt}{1.0em}{0.35em}

\setlength{\parindent}{0pt}
\setlength{\parskip}{0.55em}
\setlength{\tabcolsep}{4pt}
\renewcommand{\arraystretch}{1.18}
\setlist[itemize]{leftmargin=1.35em,itemsep=0.25em,topsep=0.25em}
\setlist[enumerate]{leftmargin=1.55em,itemsep=0.3em,topsep=0.3em}

\pagestyle{fancy}
\fancyhf{}
\fancyhead[L]{\sffamily\scriptsize COMPACTINVACUUM POLARIMETER}
\fancyhead[R]{\sffamily\scriptsize PRELIMINARY ENGINEERING REPORT}
\fancyfoot[L]{\sffamily\scriptsize Not for fabrication or purchase-order release}
\fancyfoot[R]{\sffamily\scriptsize Page \thepage}
\renewcommand{\headrulewidth}{0.25pt}
\renewcommand{\footrulewidth}{0pt}

\newcolumntype{Y}{>{\raggedright\arraybackslash}X}
\newcolumntype{C}[1]{>{\centering\arraybackslash}p{#1}}
\newcolumntype{L}[1]{>{\raggedright\arraybackslash}p{#1}}

\newtcolorbox{decisionbox}[1]{
  enhanced,
  breakable,
  colback=ReportGreen,
  colframe=ReportBlue!45,
  boxrule=0.55pt,
  arc=1.2mm,
  left=2.2mm,
  right=2.2mm,
  top=1.8mm,
  bottom=1.8mm,
  title={#1},
  fonttitle=\sffamily\bfseries,
  coltitle=ReportDarkBlue
}

\newtcolorbox{warningbox}[1]{
  enhanced,
  breakable,
  colback=ReportRed,
  colframe=red!45!black,
  boxrule=0.55pt,
  arc=1.2mm,
  left=2.2mm,
  right=2.2mm,
  top=1.8mm,
  bottom=1.8mm,
  title={#1},
  fonttitle=\sffamily\bfseries,
  coltitle=red!45!black
}

\newtcolorbox{notebox}[1]{
  enhanced,
  breakable,
  colback=ReportAmber,
  colframe=orange!55!black,
  boxrule=0.5pt,
  arc=1.2mm,
  left=2.2mm,
  right=2.2mm,
  top=1.6mm,
  bottom=1.6mm,
  title={#1},
  fonttitle=\sffamily\bfseries,
  coltitle=orange!45!black
}

\newcommand{\datatag}[1]{%
  \par{\sffamily\scriptsize\color{ReportMuted}Data origin: \path{#1}}\par
}
\newcommand{\proposal}{\textsc{Proposed}}
\newcommand{\derived}{\textsc{Derived}}
\newcommand{\sourcevalue}{\textsc{Source}}
\newcommand{\provisional}{\textsc{Provisional}}

% Generated by scripts/generate_report_data.py; do not edit.
\newcommand{\ReportGitCommit}{\detokenize{cb57cdf88ebb66067e5f42d8f10175e0d6fe3afd}}
\newcommand{\BeamEnergyMeV}{380.0}
\newcommand{\ActiveDiameterMm}{20.0}
\newcommand{\ActiveLengthMm}{5.5}
\newcommand{\ActiveThicknessMm}{5.5}
\newcommand{\LegacyThicknessMm}{10.0}
\newcommand{\CandidateMinimumDepositMeV}{1.88}
\newcommand{\CandidateMaximumDepositMeV}{15.15}
\newcommand{\LegacyMaximumDepositMeV}{35.98}
\newcommand{\MaximumDepositReductionPercent}{58}
\newcommand{\WorstCaseNormalRangeMm}{594.79}
\newcommand{\LisePrimaryRangeModel}{ATIMA 1.2 LS}
\newcommand{\ForwardCmLow}{63.26}
\newcommand{\ForwardCmHigh}{74.13}
\newcommand{\BackwardCmLow}{149.21}
\newcommand{\BackwardCmHigh}{162.22}
\newcommand{\EqrNormalNonlinearityTwo}{0.41}
\newcommand{\EqrNormalNonlinearityFive}{1.01}
\newcommand{\EqrNormalCharge}{0.216}
\newcommand{\EqrLegacyNormalNonlinearityFive}{2.38}
\newcommand{\EqrLegacyNormalCharge}{0.507}
\newcommand{\EqrCarbonNonlinearity}{11.89}
\newcommand{\FrontEndAverageCurrentMilliAmp}{10.8}
\newcommand{\FrontEndVoltageScale}{0.54}
\newcommand{\LowGainRangeNc}{1.5}
\newcommand{\ChamberInnerXmm}{440}
\newcommand{\ChamberInnerYmm}{440}
\newcommand{\ChamberBodyLengthMm}{420}
\newcommand{\ChamberWallThicknessMm}{8}
\newcommand{\ChamberMaterialVolumeL}{9.12}
\newcommand{\ChamberScreeningMassKg}{72.0}
\newcommand{\PipeRadialWallMm}{1.65}
\newcommand{\IcfAssemblyLeakMantissa}{1.0}
\newcommand{\MaintenanceAccessStandard}{ICF305}
\newcommand{\MaintenanceAccessClearBoreMm}{251.0}
\newcommand{\MaintenanceAccessPassageMarginMm}{23.8}
\newcommand{\MaintenanceExtractionStatus}{PLACEHOLDER}
\newcommand{\SourceValidationStatus}{PASS}
\newcommand{\SourceValidationMode}{prototype\_non\_strict}


\begin{document}
\hypersetup{pageanchor=false}

\begin{titlepage}
  \thispagestyle{empty}
  \vspace*{0.62in}

  {\sffamily\bfseries\fontsize{25}{29}\selectfont\color{ReportText}
  CompactInVacuum Polarimeter\par}
  \vspace{0.25em}
  {\sffamily\fontsize{14}{18}\selectfont\color{ReportMuted}
  SiPM Detector, Vacuum Chamber, and China Procurement Engineering Report\par}

  \vspace{0.55in}
  \renewcommand{\arraystretch}{1.35}
  \begin{tabularx}{\textwidth}{>{\sffamily\bfseries}L{0.23\textwidth}Y}
    \toprule
    Status & Preliminary Design Review / Procurement Planning \\
    Issue date & 26 July 2026 \\
    CAD reference & Polarimeter commit \texttt{\ReportGitCommit} \\
    Configuration & \path{compactInVacuum/config/default_compactInVacuum.yaml} \\
    Source validation & \SourceValidationStatus \\
    Procurement region & People's Republic of China preferred \\
    Primary integration candidate & Shanghai PrMat Precision Instrument Technology Co., Ltd. (conditional qualification) \\
    \bottomrule
  \end{tabularx}

  \vfill
  \begin{warningbox}{PRELIMINARY -- NOT FOR FABRICATION OR PURCHASE-ORDER RELEASE}
    This report converts the current physics and CAD state into a proposed engineering baseline.
    Open dimensions, vacuum requirements, interface drawings, material certificates, and supplier
    capabilities must be closed at the stated development gates.
  \end{warningbox}

  \vspace{0.22in}
  {\small\color{ReportMuted}
  The numerical tables in this LaTeX report are generated from version-controlled source files.
  The complete reproduction command is \texttt{make verify}; see \cref{sec:verification}.}
\end{titlepage}

\hypersetup{pageanchor=true}
\pagenumbering{roman}
\tableofcontents
\clearpage
\pagenumbering{arabic}

\section{Executive decision}

\begin{decisionbox}{Recommended first-prototype baseline}
  \begin{itemize}
    \item Use a fast blue plastic scintillator, not a high-light-yield inorganic crystal.
    \item Use a \SI{\CompactScaledActiveDiameter}{\milli\meter} diameter by
      \SI{\ActiveThicknessMm}{\milli\meter} active element as the first 0.8-scale candidate,
      and test \SIrange{5}{6}{\milli\meter} before freezing production thickness.
      If the present radii are retained, keep the current \SI{\ActiveDiameterMm}{\milli\meter}
      diameter and use the same active thickness.
    \item Use the NDL EQR15 11-6060D-S as the baseline SiPM. Keep the NDL EQR20 as a
      higher-gain backup and a Hamamatsu S13360-6025CS as an imported reference.
    \item Design and measure the total optical collection into the SiPM active area in the
      2--5\% range. Do not treat this efficiency as a vendor specification.
    \item Keep active electronics outside vacuum. Carry each fast signal through one shielded
      \SI{50}{\ohm} coaxial feedthrough, with SiPM bias supplied through an external bias tee.
  \end{itemize}
\end{decisionbox}

At the \SI{\ActiveThicknessMm}{\milli\meter} candidate's maximum calculated normal elastic
d--p deposit of \SI{\CandidateMaximumDepositMeV}{\mega\electronvolt}, the EQR15 short-pulse microcell
nonlinearity is \EqrNormalNonlinearityTwo\% at 2\% optical collection and
\EqrNormalNonlinearityFive\% at 5\%. At 5\%, the associated primary-avalanche charge is
\SI{\EqrNormalCharge}{\nano\coulomb}. The SiPM is therefore viable, but the front-end can
still clip before the microcells become unacceptable. The first prototype must provide a
timing path, a protected low-gain charge path, and a calibrated light-injection scan.
For comparison, the legacy \SI{\LegacyThicknessMm}{\milli\meter} calculation reaches
\SI{\LegacyMaximumDepositMeV}{\mega\electronvolt}, \EqrLegacyNormalNonlinearityFive\%
nonlinearity, and \SI{\EqrLegacyNormalCharge}{\nano\coulomb} at the same optical efficiency.

The calculated carbon background point is deliberately conservative: it applies the
electron-equivalent reference light yield without Birks quenching. Its EQR15 microcell
nonlinearity is \EqrCarbonNonlinearity\% at 5\% optical collection. Heavy-ion quenching
will reduce the real occupancy, but the correction must be measured rather than assumed.

Shanghai PrMat is a reasonable candidate for vacuum-chamber integration because it publicly
describes custom vacuum-component design and manufacturing capability
\cite{prmat_about,prmat_kcell}. This is candidate evidence, not prequalification. The
purchase route must require a formal RFQ, design-for-manufacture review, weld and inspection
plans, dimensional inspection, cleaning records, and helium-leak acceptance.

\begin{warningbox}{Fixed ICF114 contract; CAD bore still requires correction}
  Two polarimeters are planned. The afterSRC external-detector instrument has fixed
  ICF114 front and rear beam interfaces and a top ICF70 rotary-feedthrough interface.
  The CompactInVacuum beam interfaces are also ICF114, and its rotary-feedthrough mounting
  interface is now fixed as ICF70. The purchased rotary unit remains unfrozen. Adapters and
  bellows are permitted, but they must preserve the specified metal-seal vacuum class.

  The current CAD configuration uses a \SI{63.6}{\milli\meter} pipe outside diameter and a
  \SI{63.0}{\milli\meter} inside diameter. The calculated radial wall is only
  \SI{\PipeRadialWallMm}{\milli\meter}. A reference ICF114 half nipple has approximately
  \SI{60.2}{\milli\meter} clear bore, so the present \SI{63}{\milli\meter} stay-clear exceeds
  that reference by \SI{\IcfOneOneFourDeficitMm}{\milli\meter}
  \cite{cosmotec_icf}. ICF114 is not an open selection item. Freeze the usable beam aperture
  from the certified ICF114 flange, nipple, adapter, and bellows drawings, then correct the
  internal pipe and stay-clear geometry before releasing fabrication.
\end{warningbox}

\begin{table}[htbp]
  \centering
  \small
  \caption{Decision summary and release conditions.}
  \begin{tabularx}{\textwidth}{L{0.18\textwidth}Y Y}
    \toprule
    Item & Recommended decision & Release condition \\
    \midrule
    Scintillator & Fast blue plastic; \SIrange{5}{6}{\milli\meter} first-prototype study band &
      Source/beam timing, threshold, light-output, and thickness-comparison data \\
    SiPM & EQR15 baseline; EQR20 backup; Hamamatsu reference &
      Bench linearity, gain, dark rate, temperature coefficient, and lot screening \\
    Chamber & Short cylindrical pressure shell plus removable flat service plate or
      faceted detector cartridge & Port map, external-pressure FEA, handling mass, and access review \\
    ICF interfaces & ICF114 front/rear beam interfaces; afterSRC top rotary interface ICF70 &
      Purchased-component drawings, metal-seal assembly, and ICF-class leak acceptance \\
    PrMat workshare & Chamber, rotary-feedthrough integration, metal cassettes, final
      dimensional and leak acceptance & Accepted capability questionnaire and quality plan \\
    Detector workshare & Scintillator supplier plus project electronics team &
      A golden cassette passes optical, vacuum, electrical, and coincidence tests \\
    \bottomrule
  \end{tabularx}
\end{table}

\section{Scope, reference configuration, and decision state}

The project contains two distinct polarimeters. The afterSRC instrument places its detectors
outside a dedicated target chamber and retains the validated H+V+V mechanical family. Its
front and rear beam interfaces are fixed as ICF114, its top rotary-feedthrough interface is
ICF70, and it has no additional side vacuum ports. CompactInVacuum is the second instrument
and the main future-development line; its detectors operate inside vacuum and its current
front/rear beam-interface configuration is also ICF114. Its top rotary mounting interface is
ICF70; four sector-grouped ICF70 four-channel coax service interfaces and one ICF70
32-pin housekeeping interface are the current concept baseline. This report focuses on
CompactInVacuum detector technology while preserving the afterSRC interface contract.

The reference physics case is a \SI{\BeamEnergyMeV}{\mega\electronvolt} deuteron beam on a
CH2 target with elastic d--p coincidence in four azimuthal sectors.

\begin{table}[htbp]
  \centering
  \small
  \caption{Detector geometry read from the active CompactInVacuum configuration and
  checked against the generated channel manifest.}
  \begin{tabular}{L{0.10\textwidth}L{0.13\textwidth}C{0.12\textwidth}C{0.20\textwidth}L{0.24\textwidth}}
    \toprule
    Station & Particle & Angle (\si{\degree}) & Active-center radius (\si{\milli\meter}) &
      Lab acceptance (\si{\degree}) \\
    \midrule
    % Generated data rows; table rules and headings are owned by the report source.
D & Deuteron & 20.9 & 140.0 & 16.81--24.99 \\
P-small & Proton & 11.2 & 190.0 & 8.19--14.21 \\
P-large & Proton & 53.4 & 205.0 & 50.61--56.19 \\
\bottomrule

  \end{tabular}
  \datatag{generated/geometry_table.tex; generated/provenance.json}
\end{table}

The current CAD detector envelope is \SI{\ActiveDiameterMm}{\milli\meter} diameter by
\SI{\ActiveLengthMm}{\milli\meter}. Its configuration still records both active medium and
photosensor as \enquote{undecided/placeholder}. The SiPM and plastic-scintillator choices in
this report are therefore a \proposal{} prototype baseline. They should be promoted into
the authoritative CompactInVacuum target/configuration only after the first optical coupon
and packaging decisions are accepted.

\subsection{Current chamber envelope}

\begin{table}[htbp]
  \centering
  \small
  \caption{Current square-vessel screening metrics. The closed-shell mass estimate ignores
  ports, weld preparations, local reinforcements, and removed material.}
  \begin{tabularx}{\textwidth}{L{0.26\textwidth}L{0.23\textwidth}Y}
    \toprule
    Parameter & Current model / calculation & Engineering interpretation \\
    \midrule
    Inner cross-section & \SI{440}{\milli\meter} by \SI{440}{\milli\meter} &
      Larger than expected for a compact detector cartridge; drives transport volume \\
    Nominal wall thickness & \SI{10}{\milli\meter} &
      Must not be frozen without external-pressure buckling analysis and weld details \\
    Body length & \SI{360}{\milli\meter} &
      A short body, but detector services and target access still dominate packaging \\
    Closed-shell material volume & \SI{\ChamberMaterialVolumeL}{\liter} &
      Arithmetic screening value before subtracting ports \\
    Stainless screening mass & \SI{\ChamberScreeningMassKg}{\kilogram} &
      Confirms that the current square body is inconvenient for manual transport \\
    Detector count & 12 operating channels &
      Three stations in each of four sectors; purchase plan must include spares \\
    \bottomrule
  \end{tabularx}
  \datatag{compactInVacuum/config/default_compactInVacuum.yaml; generated/derived_results.json}
\end{table}

\subsection{Compact v2 direction}

\begin{itemize}
  \item Evaluate a 0.8 geometric scale first: D, P-small, and P-large active-center radii
    near \SI{\CompactScaledDeuteronRadius}{\milli\meter},
    \SI{\CompactScaledProtonSmallRadius}{\milli\meter}, and
    \SI{\CompactScaledProtonLargeRadius}{\milli\meter}, with a
    \SI{\CompactScaledActiveDiameter}{\milli\meter} active diameter.
  \item Evaluate a short cylindrical pressure shell around a removable flat or faceted
    detector/service cartridge. A cylinder remains structurally efficient under external
    pressure; the flat service element simplifies ICF bosses, feedthroughs, datums, and maintenance.
  \item Group electrical services on one removable plate or top/rear manifold. Avoid
    distributing twelve signal feedthroughs around the curved shell unless routing and weld
    distortion analyses justify it.
  \item Use removable beamline spools, knife-edge covers, lifting points, and a stated
    maximum manually handled subassembly mass.
\end{itemize}

\begin{notebox}{Provisional inputs that still require closure}
  Beam energy and d--p reaction are source-controlled. The \SIrange{5}{6}{\milli\meter}
  prototype thickness band, 2--5\% total optical collection,
  \SI{10000}{photons/MeV} electron-equivalent light yield, populated-detector bake regime,
  and rotary-feedthrough load case remain engineering assumptions. The ICF114/ICF70 interface
  family and ICF-class metal-seal boundary are fixed; the operating pressure and gas-load
  budget still require beamline-owner approval.
\end{notebox}

\section{Detector physics and active thickness}

The elastic coincidence acceptance is recalculated from the current detector diameter,
active-center radii, and laboratory angles. The forward and backward common center-of-mass
windows are \ForwardCmLow--\ForwardCmHigh\si{\degree} and
\BackwardCmLow--\BackwardCmHigh\si{\degree}, respectively. Energy deposition is then
calculated with the repository's relativistic two-body kinematics and C8H8 range tables.
The primary values use \LisePrimaryRangeModel; the complete CSV includes five LISE++ range
models so interpolation-model spread is visible.

\begin{table}[htbp]
  \centering
  \scriptsize
  \caption{Calculated energy deposition in a \SI{\ActiveThicknessMm}{\milli\meter}
  C8H8 active layer. Carbon values are energy deposits, not electron-equivalent light output.}
  \begin{tabular}{L{0.22\textwidth}C{0.08\textwidth}C{0.17\textwidth}C{0.20\textwidth}L{0.17\textwidth}}
    \toprule
    Reaction / particle & Station & Lab range (\si{\degree}) &
      Deposit in \SI{\ActiveThicknessMm}{\milli\meter} plastic
      (\si{\mega\electronvolt}) & Trace status \\
    \midrule
    % Generated data rows; table rules and headings are owned by the report source.
d-p forward, deuteron & D & 19.04--22.68 & 2.92--3.19 & derived; atima\_1\_2\_ls \\
d-p forward, proton & P-large & 49.91--56.89 & 3.16--4.15 & derived; atima\_1\_2\_ls \\
d-p backward, deuteron & D & 15.80--25.44 & 9.69--13.93 & derived; atima\_1\_2\_ls \\
d-p backward, proton & P-small & 7.58--14.96 & 1.71--1.76 & derived; atima\_1\_2\_ls \\
d-C elastic, deuteron & D & 15.80--26.00 & 2.42--2.46 & derived; atima\_1\_2\_ls \\
d-C elastic, carbon & P-small & 7.44--14.96 & 182.88--192.75 & derived, unquenched; atima\_1\_2\_ls \\
d-C elastic, carbon & P-large & 49.91--56.89 & 58.15--80.89 & derived, unquenched; atima\_1\_2\_ls \\
\bottomrule

  \end{tabular}
  \datatag{generated/energy_deposition.csv; code/data/energy_lise/dp_dc_kinematics.py}
\end{table}

\subsection{Active-thickness recommendation}

The detector acts as a $\Delta E$ counter, not a stopping calorimeter. Accepted d--p
particles have calculated ranges up to approximately \SI{\WorstCaseNormalRangeMm}{\milli\meter};
none of the scanned \SIrange{1}{10}{\milli\meter} options stops every branch. Thickness must
therefore be chosen from timing/threshold margin, optical uniformity, saturation, mechanics,
and coincidence separation rather than from a false full-stopping criterion.

\begin{table}[htbp]
  \centering
  \scriptsize
  \caption{LISE++ thickness scan. Each cell gives the deposited-energy range in
  \si{\mega\electronvolt} for the accepted coincidence branch using
  \LisePrimaryRangeModel. The accompanying generated CSV preserves the full five-model result.}
  \begin{tabular}{C{0.08\textwidth}C{0.16\textwidth}C{0.16\textwidth}
                  C{0.16\textwidth}C{0.16\textwidth}L{0.10\textwidth}}
    \toprule
    Thickness (\si{\milli\meter}) & Forward d & Forward p &
      Backward d & Backward p & Role \\
    \midrule
    % Generated data rows; table rules and headings are owned by the report source.
1 & 0.58--0.63 & 0.65--0.80 & 1.91--2.43 & 0.34--0.35 & scan \\
2 & 1.17--1.26 & 1.30--1.60 & 3.85--4.97 & 0.68--0.71 & scan \\
3 & 1.75--1.89 & 1.95--2.40 & 5.86--7.65 & 1.02--1.06 & scan \\
4 & 2.33--2.52 & 2.60--3.20 & 7.92--10.49 & 1.37--1.41 & scan \\
5 & 2.92--3.16 & 3.25--4.00 & 10.05--13.54 & 1.71--1.76 & scan \\
5.5 & 3.21--3.47 & 3.57--4.42 & 11.14--15.15 & 1.88--1.94 & candidate \\
6 & 3.50--3.79 & 3.90--4.83 & 12.23--16.83 & 2.05--2.12 & scan \\
8 & 4.67--5.05 & 5.23--6.50 & 16.87--24.60 & 2.73--2.82 & scan \\
10 & 5.83--6.32 & 6.57--8.17 & 21.94--35.98 & 3.42--3.53 & legacy \\
\bottomrule

  \end{tabular}
  \datatag{generated/thickness_summary.csv; generated/thickness_scan.csv;
  code/data/energy_lise/lise_range.py}
\end{table}

The \SI{\ActiveThicknessMm}{\milli\meter} candidate gives
\SIrange{\CandidateMinimumDepositMeV}{\CandidateMaximumDepositMeV}{\mega\electronvolt}
over the normal accepted branches and reduces the maximum deposited energy by approximately
\MaximumDepositReductionPercent\% relative to \SI{\LegacyThicknessMm}{\milli\meter}.
It is the compactness/low-saturation starting point, not a frozen optimum. Build at least
\SI{5}{\milli\meter} and \SI{6}{\milli\meter} otherwise-identical cassettes; retain
\SI{4}{\milli\meter}, \SI{8}{\milli\meter}, and \SI{10}{\milli\meter} coupons for threshold,
timing, quenching, and dynamic-range comparison. Freeze production thickness only after a
source or beam test and a version-controlled Geant4 response study.

\subsection{Light-yield quenching}

The saturation model deliberately uses the electron-equivalent light yield as an upper bound.
For heavy charged particles, Birks' law gives the first-order relation
\begin{equation}
  \frac{\mathrm{d}L}{\mathrm{d}x}
  =
  S\,\frac{\mathrm{d}E/\mathrm{d}x}
  {1 + k_B\,\mathrm{d}E/\mathrm{d}x}.
\end{equation}
The supplier batch and the final surface treatment can change the effective light collection,
while the Birks constant depends on material and measurement conditions. No assumed Birks
constant is applied in the saturation numbers. Carbon values are therefore intentionally
conservative for SiPM occupancy but still valuable for defining the calibration range.

\section{Scintillator and detector cassette}

\subsection{Material choice}

\begin{table}[htbp]
  \centering
  \small
  \caption{Qualitative scintillator trade study for the first compact detector.}
  \begin{tabularx}{\textwidth}{L{0.15\textwidth}Y Y L{0.17\textwidth}}
    \toprule
    Material & Main advantage & Main concern & Disposition \\
    \midrule
    Fast blue plastic & Fast coincidence timing; easy machining; low cost; domestic supply route &
      Lower energy resolution; batch and surface quality must be measured & Baseline \\
    CsI(Tl) & High light output and widespread availability &
      Slow decay, hygroscopic handling, severe SiPM saturation risk & Not recommended \\
    GAGG(Ce) & High light output; non-hygroscopic &
      Dense high-Z material, slower timing, higher cost, more saturation & Future comparison only \\
    LYSO(Ce) & Fast inorganic option; high stopping power &
      High cost, high-Z background, intrinsic activity, saturation risk & Not recommended \\
    \bottomrule
  \end{tabularx}
\end{table}

Domestic plastic material and rough/finish machining should be quoted through Shanghai EPIC
Crystal / Shanghai Shuojie or another qualified Chinese producer
\cite{epic_plastic}. EJ-200/204-class properties are a reference, not a mandatory imported
material \cite{eljen_plastic}. The purchase specification must state the measured emission
spectrum, rise/decay behavior, light output relative to a named reference, bulk attenuation,
dimensional tolerance, surface finish, reflector, adhesive, and vacuum-material evidence.

\subsection{Recommended cassette and optical variants}

\begin{itemize}
  \item Active plastic: \SI{20}{\milli\meter} diameter by
    \SIrange{5}{6}{\milli\meter} for the 0.8-scale study, or
    \SI{25}{\milli\meter} diameter by the same tested thickness at the current radii.
  \item Polish the SiPM-coupling face and define the side/rear finish on the drawing.
  \item Couple one centered 6-by-6-mm SiPM to the flat rear face. Avoid a long acrylic light
    guide until uniformity measurements show it is necessary.
  \item Use a removable cassette with a spring-preloaded optical joint. Avoid permanent
    optical epoxy in the first prototype.
  \item Provide a copper thermal path to the chamber wall, temperature sensing near the
    SiPM, strain relief, wrapped coax, and an interchangeable mechanical datum.
\end{itemize}

The first coupon program should compare diffuse reflector, specular reflector, a shallow
integral taper, and a deliberately attenuated interface. Each variant must be scanned in two
dimensions with a collimated source while recording pulse charge, timing, spatial uniformity,
and the inferred total optical collection.

\section{SiPM selection and saturation calculation}

\begin{table}[htbp]
  \centering
  \scriptsize
  \caption{Candidate SiPM input data. PDE and gain are engineering values copied into the
  auditable assumptions file; current signed datasheets govern procurement.}
  \begin{tabular}{L{0.23\textwidth}C{0.14\textwidth}C{0.08\textwidth}C{0.12\textwidth}C{0.08\textwidth}L{0.20\textwidth}}
    \toprule
    Candidate & Active area (\si{\milli\meter\squared}) & Pitch (\si{\micro\meter}) &
      Cells & PDE & Role \\
    \midrule
    % Generated data rows; table rules and headings are owned by the report source.
NDL EQR15 11-6060D-S & 6.14$\times$6.14 & 15 & 167,537 & 45.0\% & Recommended baseline \\
NDL EQR20 11-6060D-S & 6.24$\times$6.24 & 20 & 97,344 & 47.8\% & Higher-gain backup \\
Hamamatsu S13360-6025CS & 6.00$\times$6.00 & 25 & 57,600 & 25.0\% & Imported reference \\
Ray-Quant JSP-TP6050 & 6.00$\times$6.00 & 50 & 13,852 & 35.0\% & Comparison only \\
\bottomrule

  \end{tabular}
  \datatag{data/report_assumptions.yaml; references.bib}
\end{table}

The EQR15 is preferred because its \SI{15}{\micro\meter} pitch provides the largest microcell
population among the evaluated 6-by-6-mm devices. The EQR20 has similar PDE and higher gain
but fewer cells. The Hamamatsu device is a documented imported reference. The domestic
\SI{50}{\micro\meter} Ray-Quant device has too few cells for the expected direct energy
signal unless optical coupling is deliberately attenuated
\cite{centao_eqr15,centao_eqr20,hamamatsu_s13360,rayquant_jsp}.

\subsection{Model}

For a scintillation pulse shorter than the cell-recovery time and uniformly illuminating the
active area, the first-order occupancy model is
\begin{align}
  N_{\mathrm{seed}} &=
    E_{\mathrm{dep}}\,Y_{\gamma}\,\eta_{\mathrm{opt}}\,\mathrm{PDE}, \\
  N_{\mathrm{fired}} &=
    N_{\mathrm{cell}}
    \left[1-\exp\left(-\frac{N_{\mathrm{seed}}}{N_{\mathrm{cell}}}\right)\right], \\
  \delta_{\mathrm{sat}} &=
    1-\frac{N_{\mathrm{fired}}}{N_{\mathrm{seed}}}, \\
  Q &= e\,G\,N_{\mathrm{fired}}.
\end{align}
This follows the standard short-pulse occupancy picture
\cite{hamamatsu_occupancy}. Correlated avalanches, recovery during the pulse, nonuniform
illumination, temperature, and electronics are excluded. Those effects belong in measured
channel calibration.

\begin{figure}[htbp]
  \centering
  \begin{tikzpicture}
    \begin{axis}[
      width=\textwidth,
      height=0.49\textwidth,
      xmin=0,
      xmax=210,
      ymin=0,
      ymax=56,
      xlabel={Deposited energy in unquenched light upper-bound model (\si{\mega\electronvolt})},
      ylabel={Microcell nonlinearity (\%)},
      grid=major,
      axis on top,
      legend style={font=\scriptsize,at={(0.02,0.98)},anchor=north west,fill=white,draw=gray!45},
      tick label style={font=\small},
      label style={font=\small}
    ]
      \draw[fill=green!10,draw=none] (axis cs:0,0) rectangle (axis cs:210,5);
      \addplot[ReportBlue,thick] table[x=energy_mev,y=eqr15_nonlinearity_pct,col sep=comma]
        {generated/sipm_curve.csv};
      \addlegendentry{NDL EQR15}
      \addplot[green!55!black,thick] table[x=energy_mev,y=eqr20_nonlinearity_pct,col sep=comma]
        {generated/sipm_curve.csv};
      \addlegendentry{NDL EQR20}
      \addplot[violet!75!black,thick] table[x=energy_mev,y=hamamatsu_nonlinearity_pct,col sep=comma]
        {generated/sipm_curve.csv};
      \addlegendentry{Hamamatsu S13360}
      \addplot[red!75!black,thick] table[x=energy_mev,y=rayquant_nonlinearity_pct,col sep=comma]
        {generated/sipm_curve.csv};
      \addlegendentry{Ray-Quant JSP-TP6050}
      \addplot[orange!70!black,dashed,thick] coordinates {(0,10) (210,10)};
      \addplot[ReportBlue,dotted,thick]
        coordinates {(\CandidateMaximumDepositMeV,0) (\CandidateMaximumDepositMeV,56)};
      \addplot[black!65,dotted,thick]
        coordinates {(\LegacyMaximumDepositMeV,0) (\LegacyMaximumDepositMeV,56)};
      \addplot[black!65,dotted,thick] coordinates {(192.75,0) (192.75,56)};
      \node[anchor=north west,font=\scriptsize,text=ReportBlue] at (axis cs:16,18)
        {\SI{5}{\milli\meter} candidate};
      \node[anchor=north west,font=\scriptsize,text=black!70] at (axis cs:41,29)
        {\SI{10}{\milli\meter} legacy};
      \node[anchor=north east,font=\scriptsize,text=black!70,align=right] at (axis cs:190,54)
        {unquenched carbon\\upper deposit};
      \node[anchor=south west,font=\scriptsize,text=green!45!black] at (axis cs:3,1.4)
        {preferred region};
    \end{axis}
  \end{tikzpicture}
  \caption{Calculated short-pulse SiPM nonlinearity at 5\% total optical collection.}
  \datatag{generated/sipm_curve.csv}
\end{figure}

\begin{table}[htbp]
  \centering
  \scriptsize
  \caption{Calculated saturation and charge results. The carbon column is intentionally
  unquenched and therefore conservative for light output.}
  \begin{adjustbox}{width=\textwidth}
    \begin{tabular}{lccccc}
      \toprule
      Device & \multicolumn{2}{c}{\SI{\ActiveThicknessMm}{\milli\meter} candidate upper d-p deposit} &
        Charge at 5\% (\si{\nano\coulomb}) & Carbon upper bound at 5\% &
        10\% limit at 5\% (\si{\mega\electronvolt}) \\
      \cmidrule(lr){2-3}
       & Nonlinearity at 2\% & Nonlinearity at 5\% & & Nonlinearity & \\
      \midrule
      % Generated data rows; table rules and headings are owned by the report source.
NDL EQR15 11-6060D-S & 0.4\% & 1.0\% & 0.216 & 11.9\% & 159.8 \\
NDL EQR20 11-6060D-S & 0.7\% & 1.8\% & 0.456 & 20.3\% & 87.4 \\
Hamamatsu S13360-6025CS & 0.7\% & 1.6\% & 0.209 & 18.3\% & 98.9 \\
Ray-Quant JSP-TP6050 & 3.7\% & 9.0\% & 0.812 & 62.5\% & 17.0 \\
\bottomrule

    \end{tabular}
  \end{adjustbox}
  \datatag{generated/sipm_saturation.csv; generated/derived_results.json}
\end{table}

\begin{decisionbox}{Saturation decision}
  The EQR15 remains below the preferred 5\% microcell-nonlinearity target for every normal
  d--p event at 2--5\% total optical collection. The conservative unquenched carbon point is
  approximately \EqrCarbonNonlinearity\% nonlinear. The Ray-Quant \SI{50}{\micro\meter}
  comparison device is not recommended for direct coupling because its normal upper d--p
  point is already about 8.3\% nonlinear at 5\% collection.
\end{decisionbox}

\subsection{Front-end range and channel architecture}

At the \SI{\ActiveThicknessMm}{\milli\meter} candidate's normal upper d--p point and 5\%
optical collection, the EQR15 calculation gives
\SI{\EqrNormalCharge}{\nano\coulomb}. Spread over an effective
\SI{20}{\nano\second} window, this corresponds to \SI{\FrontEndAverageCurrentMilliAmp}{\milli\ampere}
average current and a \SI{\FrontEndVoltageScale}{\volt} scale across
\SI{50}{\ohm}. The real peak depends on sensor capacitance, cable, bias tee, termination,
and shaping.

\begin{itemize}
  \item Provide a fast timing-discriminator path and a lower-gain charge path, or one
    protected broadband digitizer path with two calibrated gain ranges.
  \item Set the prototype low-gain integrated-charge range to at least
    \SI{\LowGainRangeNc}{\nano\coulomb}; revise it after carbon/quenching measurements.
  \item Use programmable low-noise bias with temperature compensation. Store temperature
    and bias in run metadata.
  \item Measure breakdown voltage, gain, dark count, crosstalk, afterpulsing, recovery,
    spatial response, saturation, and temperature coefficient for every production lot.
\end{itemize}

\section{Vacuum electrical services, grounding, and thermal design}

The recommended service architecture uses one shielded \SI{50}{\ohm} coaxial path per
detector. Signal and bias share the coax through an external bias tee. Analog electronics
remain outside vacuum. For twelve detectors, the baseline is twelve separately testable
coaxial feedthrough channels plus low-rate housekeeping. The current CAD groups the channels
by azimuth sector: four ICF70 four-channel coax interfaces provide three used channels and
one installed spare per sector. A separate ICF70 32-pin envelope provides 24 used conductors
for twelve two-wire SiPM temperature sensors and eight spare pins. Exact purchased feedthrough
part numbers remain unfrozen.

Beijing Feedivac publicly lists \SI{50}{\ohm} SMA50/SMAD50 feedthrough arrangements on CF
interfaces \cite{feedivac_sma}. The RFQ must freeze exact part numbers, connector gender,
in-vacuum cable type and length, bake temperature, leak certificate, installed bandwidth,
and spare quantities. Public web data are screening evidence only.

\begin{itemize}
  \item Bond the metal cassette and chamber to protective earth; do not rely on uncertain
    mechanical contacts through removable joints.
  \item Keep the scintillator below its qualified service temperature during bake.
    Use a controlled chamber bake with the detector cartridge removable where practical.
  \item Use low-outgassing, mechanically secured internal cabling. Avoid unvented blind holes,
    virtual leaks, loose reflector film, and unsupported connectors.
  \item Monitor each SiPM temperature and define a bias-temperature correction in calibration.
\end{itemize}

\section{Chamber, ICF interfaces, and rotary feedthrough}

\subsection{Cylinder versus square chamber}

A cylindrical shell is not incompatible with ICF ports or feedthroughs. Radial ports are
normally integrated through welded bosses or formed local pads. The concern is practical:
twelve distributed signal ports, a rotary feedthrough, beam ports, and detector access can
produce a crowded curved shell with difficult datum control and maintenance.

The recommended concept is therefore hybrid:
\begin{itemize}
  \item a short cylindrical pressure shell for vacuum efficiency and lower mass;
  \item a removable flat service plate or faceted cartridge for feedthroughs, detector
    cassettes, cable routing, and alignment datums;
  \item certified ICF114 beam-port components at both ends; and
  \item a dedicated ICF70 rotary-feedthrough boss for both instruments. For
    CompactInVacuum the mounting interface is fixed, while the purchased unit, shaft-load
    rating, torque, backlash, lifetime, and bake limit remain RFQ closure items.
\end{itemize}

An initial thin-shell comparison suggests why a round vessel is attractive, but final
thickness must come from external-pressure FEA including penetrations, weld efficiencies,
local reinforcement, handling loads, and imperfection sensitivity.

\subsection{Frozen ICF interface contract}

\begin{table}[htbp]
  \centering
  \small
  \caption{Interface contract for the two polarimeters. Certified supplier drawings govern
  exact bores, bolt geometry, and installed envelope.}
  \begin{tabularx}{\textwidth}{L{0.20\textwidth}L{0.24\textwidth}L{0.22\textwidth}Y}
    \toprule
    Instrument & Beam interfaces & Other vacuum interfaces & Fixed requirement \\
    \midrule
    afterSRC & Front ICF114; rear ICF114 &
      Top ICF70 rotary feedthrough; no added side vacuum ports &
      External mating faces remain ICF114 even when adapters or bellows are used \\
    CompactInVacuum & Front ICF114; rear ICF114 &
      Top ICF70 rotary mount; four ICF70 4-channel coax interfaces; one ICF70 32-pin
      housekeeping interface &
      Mount standards and capacity are fixed; supplier models remain interface envelopes \\
    Both instruments & Purchased ICF knife edges and OFHC copper gaskets &
      Chamber, welds, adapters, bellows, and feedthroughs are one acceptance boundary &
      Assembly helium-leak acceptance $\leq
      \IcfAssemblyLeakMantissa\times10^{-10}\,\si{\pascal\meter\cubed\per\second}$ \\
    \bottomrule
  \end{tabularx}
  \datatag{docs/specs/BLP_v1_requirement_baseline.md;
  compactInVacuum/config/default_compactInVacuum.yaml; generated/derived_results.json}
\end{table}

Purchase certified ICF flanges or half nipples from a qualified vacuum-component source.
PrMat may weld and integrate them, but should not recreate the knife edge without a controlled,
qualified process. ICF uses a knife edge and OFHC copper gasket for UHV service; internal
welds and pocket-free geometry are required to avoid virtual leaks
\cite{cosmotec_icf_standard,leybold_cf_hardware}. Avoid describing every nominal 4.5-inch
CF-style product as \enquote{ICF114}; the vendor drawing and mating component must define
the bore and bolt contract.

Adapters and bellows are explicitly allowed for alignment, vibration isolation, and
installation tolerance. Each purchased item must return a signed drawing, material and
cleaning record, bake rating, and helium-leak certificate. The complete welded chamber
assembly is then leak-tested again; individual component certificates do not replace the
system test.

\subsection{Rotary-feedthrough RFQ data}

Both instruments now use a top-mounted ICF70 rotary-feedthrough mounting interface. The
current mechanism study uses an approximately \SI{18}{\milli\meter} shaft and 0/90-degree
work positions. For CompactInVacuum, fixing ICF70 does not select a supplier mechanism:
the interface envelope remains provisional until the load and service study closes.
Bellows-sealed and magnetic-coupling ICF motion-feedthrough families are commercially
available \cite{cosmotec_motion_feedthrough}. PrMat should supply or integrate the selected
unit and provide:
\begin{itemize}
  \item flange standard and complete dimensions, shaft diameter, material, usable internal
    length, external envelope, and coupling detail;
  \item allowable radial, axial, and overturning loads; output torque; breakaway torque;
    backlash; angular repeatability; hard stops; and position indication or encoder provision;
  \item seal and bearing technology, continuous versus intermittent duty, lifecycle,
    lubricant, particle generation, magnetic materials, and service/rebuild procedure; and
  \item rated operating and bake temperatures, vacuum class, certified helium-leak rate,
    motor/manual options, and interface agreement with the chamber.
\end{itemize}

\section{Chamber manufacturing and acceptance}

\begin{table}[htbp]
  \centering
  \small
  \caption{Minimum chamber manufacturing specification.}
  \begin{tabularx}{\textwidth}{L{0.18\textwidth}Y Y}
    \toprule
    Topic & Proposed RFQ requirement & Required evidence \\
    \midrule
    Material & 304L or 316L stainless steel with heat/lot traceability; final alloy by
      vacuum and magnetic review & Material certificates and traceability \\
    Purchased ICF items & Certified flanges/half nipples; no reverse-engineered knife edges &
      Supplier drawings and certificates \\
    Welding & Full weld map, process, welder qualification, sequence, and distortion plan &
      Weld plan, WPS/PQR as appropriate, weld log, visual records \\
    Vacuum design & ICF knife-edge/OFHC primary seals; no trapped volume; vented blind holes;
      cleanable internal geometry; adapters/bellows may not introduce an elastomer downgrade &
      Sectioned drawing, seal schedule, and design-review checklist \\
    Structure & External-pressure buckling plus handling and interface loads &
      Calculation report and approved drawing \\
    Dimensional control & Beam-axis datums, flange-face position/orientation, detector
      datums, and CMM plan & First-article dimensional report \\
    Leak acceptance & Complete assembly helium leak
      $\leq\IcfAssemblyLeakMantissa\times10^{-10}\,\si{\pascal\meter\cubed\per\second}$;
      record method, sensitivity, dwell, and background &
      Calibrated helium-leak report with serial number and test configuration \\
    Cleaning & Controlled process compatible with detector and vacuum materials &
      Cleaning and packing records \\
    Bake qualification & Bare-metal chamber and bought-out ICF hardware qualified separately
      from the populated detector cartridge & Component bake ratings and an approved
      assembled-system temperature limit \\
    \bottomrule
  \end{tabularx}
\end{table}

\begin{notebox}{No release before a supplier-specific package exists}
  Do not release fabrication from a STEP model alone. The fabrication package must contain
  controlled two-dimensional drawings, GD\&T, weld symbols, material and cleaning notes,
  purchased-component part numbers, acceptance criteria, and an inspection plan.
\end{notebox}

\subsection{PrMat qualification}

PrMat's public material supports candidate status, but this project requires demonstrated
vacuum precision, welded chamber experience, dimensional control, and helium-leak capability.
Before award:
\begin{itemize}
  \item request two comparable chamber references, maximum completed chamber size/mass,
    in-house vacuum capability, CMM capability, helium-leak equipment and calibration;
  \item ask whether the rotary feedthrough is a PrMat design, OEM item, or integrated
    third-party product, and obtain the manufacturer and authenticated certificate;
  \item separate supplier FEA from project-side independent review for the pressure boundary;
  \item require a process traveler covering material receipt, machining, welding, inspection,
    cleaning, assembly, and packing; and
  \item define engineering hold points after DFM, before welding, after dimensional inspection,
    and after pressure/leak acceptance.
\end{itemize}

\section{China-based procurement and workshare}

\begin{longtable}{L{0.17\textwidth}L{0.22\textwidth}L{0.31\textwidth}L{0.22\textwidth}}
  \caption{Proposed China-based workshare and qualification action.}\\
  \toprule
  Package & Preferred China route & Recommended responsibility & Qualification action \\
  \midrule
  \endfirsthead
  \toprule
  Package & Preferred China route & Recommended responsibility & Qualification action \\
  \midrule
  \endhead
  Vacuum chamber and target integration & Shanghai PrMat &
    Chamber DFM, welding, machining, assembly, dimensional and leak acceptance &
    Capability questionnaire, references, process and inspection plan \\
  Rotary feedthrough & Shanghai PrMat purchase/integration route &
    Selected model, OEM data, top ICF integration &
    Torque/load/life/leak/bake certificate and incoming inspection \\
  Plastic scintillator & Shanghai EPIC Crystal / Shanghai Shuojie or qualified domestic producer &
    Material batch, rough and precision machining, polishing, reflector options &
    Witness samples, light yield, timing, uniformity, vacuum screen \\
  Primary SiPM & NDL through Shanghai Centao or authorized domestic route &
    EQR15 11-6060D-S with current datasheet and lot traceability &
    Breakdown, gain, DCR, crosstalk, afterpulse, linearity, temperature screen \\
  Backup SiPM & NDL EQR20 through Shanghai Centao &
    Small prototype quantity &
    Same screening; compare timing and electronics headroom \\
  Reference SiPM & Hamamatsu Photonics China &
    Small control quantity &
    Reference against domestic devices \\
  Coax feedthroughs & Beijing Feedivac &
    CF40 four-channel \SI{50}{\ohm} assemblies and vacuum cables &
    Freeze part number, gender, bake, leak certificate, and spares \\
  Electronics & Chinese PCB assembler plus project design &
    Bias, protection, timing, charge range, temperature monitoring &
    Schematic review, production test, calibration and burn-in \\
  Metal cassettes & PrMat or detector-part machining supplier &
    Replaceable sensor cassette, thermal path, datums, strain relief &
    First-article CMM, fit, vacuum materials, repeated assembly test \\
  \bottomrule
\end{longtable}

PrMat should own the vacuum boundary, purchased ICF integration, rotary-feedthrough
installation, metal cassette interfaces, final chamber dimensional report, and helium-leak
acceptance. The scintillator supplier should own the optical material and polished part.
The project electronics team should own sensor selection, PCB, dynamic range, bias control,
calibration, and data interpretation. PrMat may perform final mechanical assembly only after
the project supplies an accepted golden cassette and a controlled assembly traveler.

\subsection{Preliminary purchase phasing}

\begin{table}[htbp]
  \centering
  \small
  \caption{Recommended staged purchase quantities.}
  \begin{tabularx}{\textwidth}{L{0.17\textwidth}L{0.25\textwidth}L{0.25\textwidth}Y}
    \toprule
    Gate & Suggested items & Quantity guidance & Purpose \\
    \midrule
    Gate 0: bench study & EQR15, EQR20, reference SiPMs; plastic coupons &
      At least 5 EQR15, 3 EQR20, 2 reference devices; 4 coupon styles &
      Lot screening and optical-interface choice \\
    Gate 1: one sector & Three channel cassettes, cables, feedthrough allocation &
      D, P-small, P-large plus spares &
      Coincidence, noise, timing, saturation, calibration \\
    Gate 2: chamber first article & One chamber, one rotary feedthrough, ICF components &
      One integrated first article &
      DFM, dimensional, pressure, leak, bake, maintenance tests \\
    Production & Twelve operating channels plus spares &
      At least 15 scintillators and 15 screened SiPMs &
      Matched operation with replacement inventory \\
    \bottomrule
  \end{tabularx}
\end{table}

\section{Development and verification plan}
\label{sec:verification-plan}

\begin{enumerate}
  \item \textbf{Gate 0 -- freeze physics and interfaces.} Approve beam stay-clear, detector
    radii/diameters, the \SIrange{5}{6}{\milli\meter} prototype study band, target envelope,
    populated-system bake limit, port map, survey datums, and transport limits. Record ICF114
    front/rear and afterSRC ICF70 top interfaces as fixed.
  \item \textbf{Gate 1 -- single-detector optical/electronic prototype.} Compare EQR15,
    EQR20, and the reference SiPM with 2D uniformity scans. Complete LED/laser saturation,
    charge range, timing, temperature, and 4/5/6/8/10-mm plastic-thickness comparison.
  \item \textbf{Gate 2 -- one-sector coincidence prototype.} Build D, P-small, and P-large
    channels with final-style cabling and front-end. Validate thresholds, time walk, energy
    separation, common-mode noise, and calibration metadata.
  \item \textbf{Gate 3 -- vacuum detector cartridge.} Run vacuum soak, thermal cycling,
    dark-rate monitoring, outgassing screening, cable/feedthrough noise testing, and repeated
    cassette removal/reinstallation.
  \item \textbf{Gate 4 -- engineering chamber.} Complete PrMat first article, external-pressure
    review, CMM report, helium-leak test, target-motion test, detector insertion/removal trial,
    and transport/handling demonstration.
  \item \textbf{Gate 5 -- production release.} Freeze drawings and travelers only after
    the golden sector and engineering chamber pass. Procure production channels plus approved
    spares and repeat acceptance tests.
\end{enumerate}

\begin{table}[htbp]
  \centering
  \scriptsize
  \caption{Minimum detector acceptance matrix.}
  \begin{tabularx}{\textwidth}{L{0.16\textwidth}L{0.27\textwidth}Y L{0.18\textwidth}}
    \toprule
    Test & Method & Provisional acceptance action & Produced record \\
    \midrule
    SiPM electrical screen & Automated I--V, gain, dark count, crosstalk, afterpulse,
      temperature scan & Accepted statistical distribution; outliers rejected &
      Per-device serial record \\
    Optical linearity & Pulsed 415--425-nm source with calibrated intensity steps &
      Corrected residual within project target through normal d--p range; no clipping &
      Per-channel correction curve \\
    Timing & Fast pulsed source and coincidence reference &
      Threshold and time-walk correction support final coincidence window &
      Timing calibration \\
    Uniformity & Collimated 2D source scan across active face &
      Agreed spatial uniformity; edge response documented &
      2D map and summary \\
    Vacuum/thermal & Soak and temperature cycles with dark/noise monitoring &
      Stable response; no material or connector degradation &
      Assembly serial record \\
    Mechanical repeatability & Repeated cassette removal and installation &
      Response and datum shift within project tolerance &
      Fit and calibration comparison \\
    Cable crosstalk & Source scan with neighboring channels terminated and active &
      Normal d--p pulses remain identifiable in required coincidence window &
      Configuration and crosstalk matrix \\
    \bottomrule
  \end{tabularx}
\end{table}

\section{Risk register and decisions required before purchase}

\begin{longtable}{L{0.27\textwidth}C{0.10\textwidth}C{0.10\textwidth}L{0.37\textwidth}}
  \caption{Principal risks and controls.}\\
  \toprule
  Risk & Likelihood & Impact & Control / exit condition \\
  \midrule
  \endfirsthead
  \toprule
  Risk & Likelihood & Impact & Control / exit condition \\
  \midrule
  \endhead
  True \SI{63}{\milli\meter} stay-clear is incompatible with reference ICF114 bore &
    High & High & Keep ICF114 fixed; replace placeholder pipe geometry with the certified
    nipple/adapter/bellows bore and reduce the internal stay-clear budget accordingly \\
  Scintillator thickness selected from a single LISE++ column or treated as a calorimeter &
    Medium & High & Use the five-model CSV envelope, test 5/6 mm first, retain 4/8/10 mm
    coupons, and freeze only after timing/threshold and Geant4 validation \\
  SiPM or front-end saturation on unquenched high-load events &
    Medium & High & EQR15, controlled optical collection, \SI{\LowGainRangeNc}{\nano\coulomb}
    minimum charge range, calibrated saturation correction \\
  Front-end clips before SiPM microcells saturate &
    High & High & Protected timing path, low-gain range, measured waveforms and input headroom \\
  Chamber external-pressure or weld-distortion failure &
    Low & High & Independent external-pressure review, weld sequence, dimensional hold points,
    and helium-leak acceptance \\
  Supplier capability inferred from public information &
    Medium & High & Formal capability evidence, reference projects, inspection and quality plan \\
  Optical response varies across SiPM area &
    Medium & Medium & Controlled coupling, diffuser/taper variants, 2D scan, per-channel correction \\
  Organic materials outgas or degrade in vacuum &
    Medium & High & Material evidence, witness coupon, vacuum soak, bake plan, removable cartridge \\
  \bottomrule
\end{longtable}

Before any production purchase:
\begin{itemize}
  \item approve the beam stay-clear against certified ICF114 nipple, adapter, and bellows
    drawings; ICF114 is fixed rather than a flange-size trade;
  \item approve the 0.8-scale geometry or retain current radii and \SI{25}{\milli\meter} acceptance;
  \item approve 5-mm and 6-mm active-thickness cassettes for Gate 1 and the comparison coupons;
  \item approve the ICF-class metal-seal boundary and assembly leak limit, then define operating
    pressure, gas load, populated-system bake temperature, pump-down time, cleaning class,
    and residual-gas requirement;
  \item freeze rotary-feedthrough load, torque, life, backlash, angular repeatability,
    duty cycle, and maintenance review;
  \item approve twelve-channel coax architecture, digitizer range, timing requirement,
    grounding, thermal design, and temperature/bias metadata; and
  \item accept PrMat's capability, quality plan, quotation scope, and first-article plan.
\end{itemize}

\section{RFQ checklist for Shanghai PrMat}

The RFQ should be sent as a controlled response matrix so supplier answers, attachments,
and deviations can be reviewed independently.

\begin{longtable}{L{0.19\textwidth}L{0.34\textwidth}L{0.38\textwidth}}
  \caption{Minimum RFQ response matrix for Shanghai PrMat.}\\
  \toprule
  RFQ section & Information provided by the project & Information required from PrMat \\
  \midrule
  \endfirsthead
  \toprule
  RFQ section & Information provided by the project & Information required from PrMat \\
  \midrule
  \endhead
  System overview & Approved 3D envelope, beam axis, detector zones, handling limits,
    maintenance direction & Value-to-fabricate concept, mass and COG, critical interfaces,
    maintenance and lifting provisions \\
  Vacuum boundary & Working pressure, cleaning/bake requirement, selected ICF interfaces &
    Material/weld design, venting, external-pressure calculation, leak-test method and certificate \\
  Beam interfaces & Fixed ICF114 front/rear interfaces with stay-clear and mating drawings &
    Purchased-part identity, knife-edge protection, flange position/orientation, inspection method \\
  Rotary target & Shaft and target loads, 0/90 positions, stop requirement, backlash target &
    Selected OEM/model, torque, radial/axial loads, repeatability, life, rebuild, bake, leak certificate \\
  Detector cartridge & Golden cassette interface, cable routing, thermal-path requirement &
    Detailed metal cassette, repeatable datums, cable strain relief, removal procedure \\
  Electrical feedthroughs & Twelve \SI{50}{\ohm} coax concept and housekeeping requirement &
    Selected part numbers, connector/cable drawings, certificates, spare recommendation \\
  Metrology & Critical datums, flange runout, detector axes, target-axis requirement &
    CMM plan and final report, inspection uncertainty, traceable instruments \\
  Acceptance & Vacuum soak, helium-leak and dimensional criteria &
    Inspection and test plan, calibrated equipment, nonconformance procedure \\
  Commercial & Prototype and production quantity, delivery target, change-control expectation &
    Itemized quotation, lead time, warranty, documentation list, IP/drawing ownership,
    packing and transport plan \\
  \bottomrule
\end{longtable}

Recommended order:
\begin{enumerate}
  \item freeze the rotary-feedthrough data review and one detector-cassette interface;
  \item send chamber requirements and the separate response matrix;
  \item obtain a manufacturing review before any hard tooling or welding;
  \item release production only after the chamber first article and one-sector detector
    prototype pass their acceptance gates.
\end{enumerate}

\section{Data and calculation verification}
\label{sec:verification}

The report separates four classes of information:
\begin{table}[htbp]
  \centering
  \small
  \caption{Traceability classes used throughout the report.}
  \begin{tabularx}{\textwidth}{L{0.17\textwidth}Y Y}
    \toprule
    Class & Meaning & Verification route \\
    \midrule
    \sourcevalue & Read directly from CAD configuration, channel manifest, validation report,
      or cited supplier evidence & Source hash plus direct file inspection \\
    \derived & Recalculated from source values by the report generator &
      CSV/JSON output, unit tests, and stale-file check \\
    \provisional & Engineering screening input not yet measured or frozen &
      Explicit entry in \path{data/report_assumptions.yaml} \\
    \proposal & Recommended choice not yet promoted into authoritative CAD configuration &
      Design review and controlled configuration update \\
    \bottomrule
  \end{tabularx}
\end{table}

\subsection{One-command reproduction}

From this report directory:
\begin{Verbatim}[fontsize=\small,frame=single]
make verify
\end{Verbatim}

The command performs all of the following:
\begin{enumerate}
  \item loads the current CompactInVacuum YAML configuration, target, channel manifest,
    and strict validation report;
  \item checks detector angles and active-center radii against the generated CAD manifest;
  \item reuses \path{code/data/energy_lise/dp_dc_kinematics.py} for kinematics and
    \path{code/data/energy_lise/lise_range.py} for explicit LISE++ model selection, then
    recalculates the 1/2/3/4/5/6/8/10-mm energy-deposition scan;
  \item recalculates every SiPM saturation, charge, and 10\%-nonlinearity result;
  \item recalculates chamber screening volume/mass and ICF arithmetic;
  \item compares regenerated CSV, JSON, and LaTeX fragments byte-for-byte against the files
    included by this report;
  \item runs the independent unit tests; and
  \item compiles the report with LuaLaTeX/Biber and verifies searchable PDF content.
\end{enumerate}

Useful partial commands are:
\begin{Verbatim}[fontsize=\small,frame=single]
make data
make check-data
make test
make report
\end{Verbatim}

\subsection{Generated, inspectable data}

\begin{itemize}
  \item \path{generated/derived_results.json}: full-precision derived geometry, energy,
    saturation, electronics, chamber, and ICF results.
  \item \path{generated/energy_deposition.csv}: machine-readable acceptance and energy ranges.
  \item \path{generated/thickness_summary.csv}: compact one-row-per-thickness comparison.
  \item \path{generated/thickness_scan.csv}: full event-by-event, five-model LISE++ scan.
  \item \path{generated/sipm_saturation.csv}: one-row-per-device saturation and charge summary.
  \item \path{generated/sipm_curve.csv}: all plotted saturation points.
  \item \path{generated/provenance.json}: source paths, SHA-256 hashes, Git commit, environment,
    CAD decision state, and report proposal state.
  \item \path{generated/verification_summary.txt}: human-readable PASS/FAIL record.
\end{itemize}

\begin{longtable}{L{0.18\textwidth}L{0.47\textwidth}L{0.27\textwidth}}
  \caption{Source-file provenance captured when numerical fragments were generated.}\\
  \toprule
  Source role & Repository-relative file & SHA-256 \\
  \midrule
  \endfirsthead
  \toprule
  Source role & Repository-relative file & SHA-256 \\
  \midrule
  \endhead
  % Generated data rows; table rules and headings are owned by the report source.
compact architecture baseline & \path{docs/specs/compact_one_requirement_baseline.md} & \texttt{651efe396cf3\ldots} \\
compact platform common & \path{compactInVacuum/config/common_detector.yaml} & \texttt{0b6e6552ff01\ldots} \\
compact after src profile & \path{compactInVacuum/config/afterSRC_compact.yaml} & \texttt{fe7258bfe6dd\ldots} \\
compact pre samurai profile & \path{compactInVacuum/config/infrontSamurai_compact.yaml} & \texttt{c002d966a31d\ldots} \\
compact config & \path{compactInVacuum/config/default_compactInVacuum.yaml} & \texttt{099c4f559653\ldots} \\
compact target & \path{compactInVacuum/target.yaml} & \texttt{abe34b0fbc06\ldots} \\
channel manifest & \path{compactInVacuum/compactInVacuum.channel_manifest.json} & \texttt{799760149014\ldots} \\
validation report & \path{compactInVacuum/compactInVacuum.validation_report.json} & \texttt{0701e2bebd1b\ldots} \\
analysis scenario & \path{code/config/compact_in_vacuum.ini} & \texttt{be0bc270a47a\ldots} \\
energy model & \path{code/data/energy_lise/dp_dc_kinematics.py} & \texttt{4c650dc8bd1a\ldots} \\
lise range model & \path{code/data/energy_lise/lise_range.py} & \texttt{ef329f25f9c1\ldots} \\
mechanical baseline & \path{docs/specs/BLP_v1_requirement_baseline.md} & \texttt{f4e6e1ce5980\ldots} \\
report source en & \path{docs/polarimeter_detector/compact_in_vacuum_sipm_report/compact_in_vacuum_sipm_report.tex} & \texttt{0a9757648615\ldots} \\
report source zh & \path{docs/polarimeter_detector/compact_in_vacuum_sipm_report/compact_in_vacuum_sipm_report_zh.tex} & \texttt{f3b236632754\ldots} \\
report generator & \path{docs/polarimeter_detector/compact_in_vacuum_sipm_report/scripts/generate_report_data.py} & \texttt{1abf86dc11c9\ldots} \\
report makefile & \path{docs/polarimeter_detector/compact_in_vacuum_sipm_report/Makefile} & \texttt{c7d5160535a9\ldots} \\
report assumptions & \path{docs/polarimeter_detector/compact_in_vacuum_sipm_report/data/report_assumptions.yaml} & \texttt{d0ae41aff893\ldots} \\
\bottomrule

\end{longtable}

\begin{notebox}{Configuration-state warning}
  The current CAD source still says \enquote{active medium: placeholder/undecided} and
  \enquote{photosensor: placeholder/undecided}. This report deliberately records the plastic
  scintillator and EQR15 as a proposed prototype baseline. The verification tool reports both
  states so the recommendation cannot be mistaken for an already-frozen CAD requirement.
\end{notebox}

\subsection{Limits of the current verification}

The d--p/d--C energy values use the repository's analytic two-body and LISE++ range-table
workflow. They use a continuous-slowing-down range subtraction and do not include range
straggling, wrapping/dead layers, incidence-angle path variation, Birks quenching, or detector
resolution. They are engineering screening values, not a substitute for the version-controlled
Geant4 model required before production. The SiPM model is first-order and assumes spatially uniform,
short-pulse illumination. The chamber mass is a closed-square-shell arithmetic estimate, not
an FEA model or fabrication mass. Supplier webpages support candidate selection only; signed
drawings, current datasheets, quotations, lot traceability, and certificates govern purchase.

\clearpage
\nocite{polarimeter_repo}
\printbibliography[title={References and supplier evidence}]

\vfill
\begin{center}
  {\sffamily\small\bfseries\color{ReportMuted} END OF PRELIMINARY REPORT}
\end{center}

\end{document}
